\section{The two-floor instance and the asymmetric audit}\label{si:instances}

\paragraph{Two-floor instance.} States $\{0,1\}$ and actions $u\in\{\tfr25,\tfr35\}$; state $0$ stays at $0$ under $u=\tfr25$ and exits under $u=\tfr35$, state $1$ stays at $1$ under $u=\tfr35$ and exits under $u=\tfr25$; a single uninformative observation. The masked backup gives $\Gamma_1=\{(1,0),(0,1)\}$ and, since each surviving state is kept alive indefinitely by its own action, $\Gamma_k=\Gamma_1$ and $V_k(b)=\max(b_0,b_1)$ for every $k\ge1$ (verified to $k=5$). The one-step exit probabilities are $p(0,\tfr35)=p(1,\tfr25)=1$, so at $b=(\tfr12,\tfr12)$ the chance-constrained certificate $\sum_xb(x)p(x,a)\ge\tfr12$ fires for both actions, giving $\Delta_k\ge\tfr12$, which is attained. The witness vectors live on $\calV\cup\{\bot\}$ as $(1,0,0)$ and $(0,1,0)$; the suppressed bottom coordinate is identically zero.

\paragraph{Asymmetric audit.} On the delayed instance at $(z_0,\Tobs)=(\tfr32,2)$ with prior mass $w=\tfr3{10}$ on $\theta=+1$, Table~\ref{tab:asym} lists the four open-loop two-step segments. The first move sends one branch to $\tfr12<1$, so every segment loses a branch at step $1$; the best segments save the heavier branch and $V_2(b_0)=\tfr7{10}$, deficit $\tfr3{10}=\min(w,1-w)$.

\begin{table}[h]
\centering\footnotesize
\caption{The asymmetric audit at $(z_0,\Tobs)=(\tfr32,2)$, $b_0=(w,1-w)=(\tfr3{10},\tfr7{10})$.}
\label{tab:asym}
\begin{tabular}{@{}l c c l@{}}
\toprule
Segment $(u_0,u_1)$ & positions after step 1 ($\theta=+1$, $\theta=-1$) & surviving mass & lost branch\\
\midrule
$(+1,+1)$ & $5/2$, $1/2$ & $\tfrac{3}{10}$ & $\theta=-1$ (mass $\tfrac{7}{10}$), at step 1\\
$(+1,-1)$ & $5/2$, $1/2$ & $\tfrac{3}{10}$ & $\theta=-1$ (mass $\tfrac{7}{10}$), at step 1\\
$(-1,+1)$ & $1/2$, $5/2$ & $\tfrac{7}{10}$ & $\theta=+1$ (mass $\tfrac{3}{10}$), at step 1\\
$(-1,-1)$ & $1/2$, $5/2$ & $\tfrac{7}{10}$ & $\theta=+1$ (mass $\tfrac{3}{10}$), at step 1\\
\bottomrule
\end{tabular}

\end{table}
